\chapter{شفرة مورس: ما نعرفه عن اللغة التي تتكلم بها عصبونات \nt{GLUT} و\nt{GABA}}
\chaptermark{شفرة مورس}
\label{sec:code}

\section{أبجدية من رمز واحد}

في شفرة مورس رمزان، النقطة والشَّرطة، وفواصل بثلاثة أطوال بينهما. أما أبجدية العصبون، كما رأينا في الفصل~\ref{sec:neurontypes}، فأفقر: رمز واحد. تدوم النبضة نحو جزء من الألف من الثانية، وكل نبضات العصبون الواحد متساوية في الارتفاع والشكل. لذلك فالرسالة كلها في الفواصل، أي في متتالية اللحظات $t_1,t_2,t_3,\dots$ إنها شفرة بلا شُّرَط، لا يحمل المعنى فيها إلا إيقاع النقاط (الشكل~\ref{fig:morse}).

\begin{figure}[t]
\centering
\begin{tikzpicture}[>=Stealth,x=0.08cm,y=1cm]
% Morse letter: dot, dash, dot
\node[left,font=\small] at (-6,2.55) {مورس:};
\fill[black] (0,2.45) rectangle (6,2.65);
\fill[black] (14,2.45) rectangle (32,2.65);
\fill[black] (40,2.45) rectangle (46,2.65);
\node[right,font=\small,gray!40!black] at (52,2.55) {نقطة، شَرطة، نقطة: المعنى في المُدد والفواصل};
% stimulus onset
\node[left,font=\small] at (-6,0.4) {العصبون:};
\draw[->,thick,green!45!black] (0,-0.55) -- (0,-0.05);
\node[below,font=\scriptsize,green!45!black] at (0,-0.55) {المنبه};
\draw[gray!70!black] (0,0) -- (152,0) node[right,black] {\small $t$};
% spikes: single ones blue, the burst red
\foreach \s in {14,26,41,88,112,131}{\draw[very thick,blue!60!black] (\s,0) -- (\s,0.8);}
\foreach \s in {60,63,66,69}{\draw[very thick,red!70!black] (\s,0) -- (\s,0.8);}
% latency
\draw[<->,thick] (0,1.1) -- node[above,font=\scriptsize] {$t_1$: زمن النبضة الأولى} (14,1.1);
% burst
\draw[decorate,decoration={brace,amplitude=4pt},red!70!black] (58,0.95) -- (71,0.95) node[midway,above=4pt,font=\scriptsize] {الرشقة هي ”الشَّرطة“};
% counting windows
\foreach \a/\b/\n in {0/50/3,50/100/5,100/150/2}{
  \draw[decorate,decoration={brace,amplitude=4pt,mirror},gray!50!black] (\a+1,-0.2) -- (\b-1,-0.2) node[midway,below=4pt,font=\small] {$n=\n$};}
\node[font=\scriptsize,gray!40!black] at (75,-1.05) {ترميز التردد: عدد النبضات $n$ في كل نافذة};
\foreach \x in {0,50,100,150}{\draw[gray!60!black] (\x,-0.06) -- (\x,0.06);}
\end{tikzpicture}
\caption{شفرة مورس وشفرة العصبون. يمكن قراءة متتالية النبضات نفسها بطرق مختلفة: عدّ النبضات في نوافذ من $50\ms$ (القسم~\ref{sec:rate})، أو قياس التأخير $t_1$~\eqref{eq:latency}، أو ملاحظة الرشقة، أي ”الشَّرطة“~\eqref{eq:burst}.}
\label{fig:morse}
\end{figure}

ما الفواصل الممكنة أصلًا؟ من الأسفل يحدها عدم الاستجابة: بعد النبضة لا يستطيع العصبون أن ينطلق ثانية مدة جزء من الألف من الثانية تقريبًا، فلا يتجاوز العدد بضع مئات من النبضات في الثانية. ومن الأعلى لا يحدها شيء: قد يصمت العصبون دقائق. ترددات عصبونات \nt{GLUT} في القشرة تتراوح بين أجزاء من النبضة وبضع عشرات من النبضات في الثانية، ومعظم العصبونات تعمل معظم الوقت قرب الحد الأدنى.

في الفصلين~\ref{sec:glutcomputer} و\ref{sec:gaba} اعتمدنا عرضًا حينًا طريقة وحينًا أخرى لكتابة الأعداد بالنبضات. والآن نسأل مباشرة: بأي لغة تتخاطب عصبونات \nt{GLUT} و\nt{GABA}؟ لا جواب كاملًا؛ فهذه اللغة لم تُفكّ رموزها. لكن بعض ما يُعرف عنها ثابت، ويمكن الوصول إلى معظمه بتفكير مهندس الاتصالات: سعة القناة، والضجيج، والمستقبِل، وكلفة الإرسال.

\section{كم يمكن أن ننقل}

لنبدأ بالحد الأعلى. ليكن المستقبِل يميز لحظات النبضات بدقة $\Delta t$: الإزاحات الأصغر من $\Delta t$ لا يميزها. لنقطّع الزمن إلى خانات طولها $\Delta t$، أصغر من زمن عدم الاستجابة، بحيث تحوي كل خانة نبضة واحدة أو لا شيء (الشكل~\ref{fig:cells}). عند تردد متوسط $r$ تقع النبضة في الخانة باحتمال $p=r\Delta t$. يحمل التدفق أكبر قدر من المعلومات حين تكون الخانات مستقلة، وعندئذ تعطي كل خانة إنتروبيا $-p\log_2p-(1-p)\log_2(1-p)\approx p\log_2(e/p)$ بت عند $p\ll1$. وبالقسمة على $\Delta t$ نحصل على السرعة القصوى للنقل وحصة النبضة الواحدة منها~\cite{MacKay1952}:
\begin{empheq}[box=\keybox]{equation}
R_{\max} \;\approx\; r\,\log_2\frac{e}{r\,\Delta t}\ \ \text{bit/s},
\qquad
\frac{R_{\max}}{r} \;\approx\; \log_2\frac{e}{r\,\Delta t}\ \ \text{bit/spike}.
\label{eq:spikeentropy}
\end{empheq}
عند $r=10$ نبضات في الثانية و$\Delta t=1\ms$ هذا نحو ثمانية بتات للنبضة وثمانين بتًا في الثانية.

\begin{figure}[t]
\centering
\begin{tikzpicture}[>=Stealth,x=0.5cm,y=1cm]
% time cut into cells of width Delta t; a cell holds one impulse or none
\foreach \k in {0,...,23}{\draw[gray!55] (\k,0) rectangle (\k+1,0.7);}
\foreach \k in {2,7,8,15,21}{
  \fill[red!10] (\k,0) rectangle (\k+1,0.7);
  \draw[very thick,red!70!black] (\k+0.5,0.05) -- (\k+0.5,0.65);}
\foreach \k in {0,...,23}{
  \pgfmathtruncatemacro{\bit}{(\k==2)||(\k==7)||(\k==8)||(\k==15)||(\k==21)}
  \ifnum\bit=1 \def\bitcol{red!70!black}\else \def\bitcol{gray!60!black}\fi
  \node[font=\scriptsize,text=\bitcol] at (\k+0.5,-0.25) {\bit};}
\draw[<->,gray!60!black] (0,0.95) -- (1,0.95) node[midway,above,font=\scriptsize,black] {$\Delta t$};
\node[font=\small,anchor=east] at (-0.3,0.35) {النبضات};
\node[font=\small,anchor=east] at (-0.3,-0.25) {البتات};
\draw[decorate,decoration={brace,amplitude=5pt,mirror},gray!60!black] (0,-0.55) -- (24,-0.55)
  node[midway,below=6pt,font=\scriptsize,black,align=center] {المستقبِل الذي يعدّ فقط: ”$n=5$“؛\\ والمستقبِل الذي يرى الخانات: أي $5$ من $24$ بالضبط، أي $\log_2\binom{24}{5}\approx15$ بتًا};
\end{tikzpicture}
\caption{من أين تأتي السعة~\eqref{eq:spikeentropy}. قُطّع الزمن إلى خانات طولها $\Delta t$، وفي كل خانة إما نبضة ($1$) أو لا ($0$): تدفق النبضات سلسلة ثنائية. تقع النبضة في الخانة باحتمال $p=r\Delta t$. كلما صغرت الخانات طالت السلسلة وكثرت البتات؛ وعدّاد النبضات يفقد تقريبًا كل ما هو مكتوب في \emph{مواضع} الآحاد.}
\label{fig:cells}
\end{figure}

البتات لكل نبضة لا تتعلق بدقة الزمن إلا لوغاريتميًا: بدقة $10\ms$ يبقى نحو خمسة بتات للنبضة. لكن إذا كان المستقبِل يعدّ النبضات في الثانية فقط، فعند $r=10$ لن يحصل على أكثر من بتّين في الثانية كلها (القسم~\ref{sec:rate}).

\begin{keyblock}
\begin{proposition}[سعة القناة النبضية]
\label{prop:capacity}
تدفق النبضات بتردد متوسط $r$، مقروءًا بدقة زمنية $\Delta t$، لا يحمل أكثر من~\eqref{eq:spikeentropy}. وتكاد السعة كلها تقبع في \emph{توقيت} النبضات: المستقبِل الذي يعدّ النبضات فقط يستخرج من التدفق نفسه أقل بعشرات المرات مما يستخرجه من يميز لحظاتها بدقة جزء من الألف من الثانية.
\end{proposition}
\end{keyblock}

هذا حد أعلى لا قياس. القياسات على العصبونات الحسية، التي يُعرف منبهها، تعطي من بت واحد إلى بضعة بتات للنبضة؛ وفي أفضل الحالات يبلغ هذا نصف الحد المحسوب لدقتها~\cite{Rieke1997}. إذن، على الأقل عند مدخل الجهاز العصبي، يُستعمل توقيت النبضات ولا يضيع.

\section{القناة الضجيجية}

حُسبت السعة~\eqref{eq:spikeentropy} بلا ضجيج. أما أين ينشأ الضجيج، فثلاث حقائق ثابتة.

\emph{مولّد النبضات نفسه دقيق.} إذا حُقن في عصبون قشري، مأخوذ من الدماغ مع قطعة من النسيج، التيارُ نفسه المتغير بسرعة في الزمن مرات كثيرة، تكررت النبضات بدقة نحو جزء من الألف من الثانية~\cite{Mainen1995}. وإذا كان التيار ثابتًا، تتبعثر لحظات النبضات من مرة إلى أخرى: فالتوقيت الدقيق تولّده التغيرات السريعة في المدخل، لا العصبون نفسه.

\emph{المشبك غير موثوق.} النبضة التي تبلغ مشبك \nt{GLUT} تقذف دفعة من الناقل العصبي باحتمال ما $p$ فقط. في مشابك الحُصين لا يصل أكثر من نصف النبضات إلى المستقبِل أصلًا، والسبب هو عشوائية القذف بالذات، لا أعطال التوصيل~\cite{AllenStevens1994}. وبالنسبة لمهندس الاتصالات هذه قناة محو؛ وعدة مشابك بين عصبونين تنقذ الموقف جزئيًا.

\emph{في القشرة الحية يبدو تدفق النبضات عشوائيًا.} الفترات بين النبضات مبعثرة تقريبًا كما في تدفق بواسون العشوائي، وعند تكرار المنبه نفسه يتغير عدد النبضات في النافذة بحيث يقارب تباينه متوسطه~\cite{Softky1993}.

كيف يولّد عنصر دقيق تدفقًا عشوائيًا؟ يتلقى العصبون آلاف المداخل، كل منها غير موثوق، والإثارة والتثبيط متوازنان كما في الفصل~\ref{sec:gaba}. يرتعش جهد الغشاء قرب العتبة، وهذه الارتعاشات هي التي تحدد لحظات النبضات. أهو ضجيج أم رسالة لا نحسن قراءتها؟ هذا إلى حد بعيد سؤال مفتوح. وقد تبيّن أن جزءًا من ”الضجيج“ رسالة فعلًا: ففي القشرة البصرية تتبع حصة كبيرة من التشتت حركات الحيوان نفسه~\cite{Stringer2019}. والصعوبة هنا مبدئية: الرسالة المضغوطة جيدًا لا تتميز عن تدفق عشوائي عند من لا يعرف الشيفرة.

\section{التردد: بطيء لكنه موثوق}
\label{sec:rate}

أبسط الشيفرات \emph{ترميز التردد}: العدد مكتوب في عدد النبضات التي يطلقها العصبون في النافذة $T$. لا المحو ولا ارتعاش اللحظات يضر بهذه الشيفرة. لكن لها ثمنًا. إذا كان التدفق بواسونيًا فعدد النبضات $n$ في النافذة متوسطه $rT$ وتشتته $\sqrt{rT}$:
\begin{empheq}[box=\keybox]{equation}
\frac{\sigma_n}{\bar n} \;=\; \frac{1}{\sqrt{r\,T}} .
\label{eq:rateerror}
\end{empheq}
لمعرفة التردد بدقة عشرة في المئة تلزم مئة نبضة، أي خمس ثوانٍ عند عشرين نبضة في الثانية. والمستويات المتجاورة تتمايز إذا تباعدت بنحو تشتتين، لذا يتمايز حتى التردد $r$ في الزمن $T$ نحو $\sqrt{rT}$ مستوى: عند $r=10$ و$T=1\,\text{s}$ هذا ثلاثة أو أربعة مستويات، أقل من بتين.
والدماغ لا يعمل بهذا البطء. يتعرف الإنسان على حيوان في صورة تُعرض عليه لحظة، وتظهر الاستجابة في الدماغ بعد نحو $150\ms$~\cite{Thorpe1996}. وبتقدير خشن تمر الإشارة في هذا الزمن بنحو عشر مراحل معالجة متتالية، بمعدل $10\text{--}20\ms$ للمرحلة. وبالترددات المعتادة في القشرة لا يلحق كل عصبون أن يطلق في المرحلة إلا صفرًا أو نبضة واحدة، فتردده في نافذة كهذه غير معرّف. وثمة مخرجان: أخذ عصبونات كثيرة، أو النظر إلى الزمن.

\emph{عصبونات كثيرة.} لتكن $N$ عصبونات تحمل العدد نفسه، والمستقبِل يجمع نبضاتها. إذا كانت ضجيجاتها مستقلة حلّ في~\eqref{eq:rateerror} $NrT$ محل $rT$: مئة عصبون عند $r=20$ في $15\ms$ تعطي ثلاثين نبضة ودقة نحو عشرين في المئة. المكان يحل محل الزمن. غير أن ضجيجات العصبونات المتجاورة في القشرة ليست مستقلة: معامل الارتباط بين أعداد النبضات لزوج من الجيران نحو $0.1$ عادة~\cite{Zohary1994}. فإذا كان يساوي $c$، فتباين المتوسط على $N$ عصبونًا يساوي
\begin{empheq}[box=\keybox]{equation}
\sigma^2_N \;=\; \sigma^2_1\,\frac{1+(N-1)\,c}{N}
\;\xrightarrow[N\to\infty]{}\; c\,\sigma^2_1 .
\label{eq:correlated}
\end{empheq}

\begin{keyblock}
\begin{proposition}[حد التوسيط]
\label{prop:pooling}
مع الضجيج المترابط لا يقلل التوسيط على مجموعة الخطأ بأكثر من $1/\sqrt{c}$ مرة. عند $c=0.1$ هذا ثلاث مرات، ويُبلغ هذا الربح ببضع عشرات من العصبونات؛ وإضافة المزيد لا تجدي.
\end{proposition}
\end{keyblock}

ليس كل ارتباط ضارًا بالقدر نفسه. لا يحد الدقة إلا ذلك الجزء من الضجيج المشترك الذي \emph{يشبه الإشارة نفسها}، أي يزيح المجموعة كلها كما كان سيزيحها تغيّر الكمية المقيسة~\cite{MorenoBote2014}. وكم من هذا الضجيج في الدماغ فعلًا، ما زال قيد البحث.

وثمة طريقة أخرى لاستعمال عصبونات كثيرة: كتابة العدد في \emph{أي} عصبون انطلق، كما في محدد اتجاه الصوت (الشكل~\ref{fig:itd}). هذا ترميز \emph{بالموضع}، وهو أوثق المعروف: في العصبونات الحسية يدل رقم السلك على أي نقطة من الجلد لُمست أو على طبقة الصوت، وهذا ثابت مرارًا. وهو مكلف في عدد العصبونات، لكنه سريع: تستغرق الرسالة تأخيرًا واحدًا.

\section{الزمن: سريع، لكن من أين نبدأ العدّ}

المخرج الثاني كتابة العدد في زمن النبضة. لنأخذ العصبون~\eqref{eq:glutneuron} ونقدّم إليه من اللحظة $t=0$ مدخلًا ثابتًا كان سيرفع الجهد إلى المستوى $U$. ينمو الجهد كـ$V=U\bigl(1-e^{-t/\tau}\bigr)$ ويبلغ العتبة $\theta$ في اللحظة
\begin{empheq}[box=\keybox]{equation}
t_1 \;=\; \tau\,\ln\frac{U}{U-\theta}, \qquad U>\theta .
\label{eq:latency}
\end{empheq}
المدخل القوي نبضة مبكرة، والضعيف نبضة متأخرة، وما دون العتبة لا نبضة. عند $\tau=10\ms$ يعطي المدخل $U=4\theta$ النبضة الأولى بعد $2.9\ms$، و$U=2\theta$ بعد $6.9\ms$، و$U=1.2\theta$ بعد $17.9\ms$. نبضة واحدة وحيدة تحمل عددًا.

لكن من أي لحظة نعدّ $t_1$؟ لا ساعة، والمستقبِل لا يعرف متى بدأ المنبه. ثمة ثلاثة أجوبة.

\emph{التأخير النسبي.} عصبونان يريان المنبه نفسه، وفرق تأخيريهما $t_1^A-t_1^B$ لا يتوقف إلا على مدخليهما، لا على لحظة البدء. وتقارن اللحظاتِ كواشفُ التزامن ذات التأخيرات (الشكل~\ref{fig:itd}). إذا أطلق $N$ عصبونًا نبضة واحدة لكل منها، فإن \emph{ترتيب} وصولها وحده يحمل حتى $\log_2N!$ بت: لعشرة عصبونات نحو اثنين وعشرين بتًا، بينما الجواب ”انطلق أم لا“ لا يعطي أكثر من عشرة. وقد بُيّن في الشبكية أن الصورة يمكن استعادتها بسرعة وموثوقية من التأخيرات النسبية للنبضات الأولى لعصبوناتها المخرجة~\cite{Gollisch2008}.

\emph{الرشقة الشبكية بداية للكلمة.} التغير الحاد في المنبه يجعل عددًا كبيرًا من العصبونات ينطلق في آن واحد تقريبًا. هذه الرشقة الشبكية نفسها علامة على البداية، كالفاصل الطويل بين الكلمات في شفرة مورس، ويعدّ المستقبِل التأخيرات منها.

\emph{الإيقاع المحلي.} في الفصل~\ref{sec:gaba} ولّدت حلقة \foreignlanguage{english}{\nt{GLUT}--\nt{GABA}} إيقاعًا محليًا دوره $12\text{--}50\ms$. ليس ساعة، لكن داخل دورته يلحق العصبون ذو المدخل القوي أن ينطلق قبل العصبون ذي المدخل الضعيف، فيتحول~\eqref{eq:latency} إلى ترميز \emph{بالطور}: العدد مكتوب في الجزء من الدورة الذي وصلت فيه النبضة. وقد رُصد هذا الترميز: العصبونات المسؤولة عن موضع معين في المكان تنطلق أبكر فأبكر في دورة إيقاع محلي بطيء كلما عبر الحيوان موضعها ”الخاص“~\cite{OKeefe1993}. أما إلى أي حد يستعمل الدماغ الطور، فغير معروف بعد.

\section{المستقبِل يختار الشيفرة}

في متتالية النبضات تردد، وأزمنة، وترتيب. أيها يُقرأ يقرره المستقبِل، ويقرره بوسيط واحد: ثابته الزمني $\tau$.

لنأخذ متوسط المعادلة~\eqref{eq:glutneuron} على الزمن. إذا وردت على العصبون تدفقات مستقلة بترددات $r_j$، فالجهد المتوسط وارتعاشه النسبي يساويان
\begin{empheq}[box=\keybox]{equation}
\bar V \;=\; \tau\sum_j w_j\,r_j ,
\qquad
\frac{\sigma_V}{\bar V} \;\approx\; \frac{1}{\sqrt{2\,r_{\text{in}}\,\tau}} ,
\label{eq:reader}
\end{empheq}
حيث المساواة الثانية مكتوبة لأوزان متساوية، و$r_{\text{in}}$ مجموع ترددات كل المداخل. حين يكون $\tau$ كبيرًا، لا يكاد الجهد يرتعش ويتبع المجموع الموزون للترددات: المستقبِل \emph{مكامل}، يقرأ الترددات وينسى الأزمنة. وحين يكون $\tau$ صغيرًا، يلحق الجهد أن يهبط بين النبضات، ولا تُبلغ العتبة إلا إذا وصلت عدة نبضات ضمن النافذة~\eqref{eq:window}: المستقبِل \emph{كاشف تزامن}، يقرأ الأزمنة (الشكل~\ref{fig:reader}). ألف مدخل بخمس نبضات في الثانية لكل منها عند $\tau=10\ms$ تعطي ارتعاشًا نحو عشرة في المئة، أي تكاملًا شبه خالص. وإذا قصّر التثبيط، كما في الفصل~\ref{sec:gaba}، $\tau$ إلى $2\ms$، نما الارتعاش إلى أكثر من عشرين في المئة، وبدأ العصبون يلاحظ التزامنات.

\begin{figure}[t]
\centering
\begin{tikzpicture}[>=Stealth,x=0.115cm,y=1cm]
% common input: the same spike train for both receivers
\foreach \s in {6,21,22,38,52,55.5,59,62.5,66,69.5,73,76.5,80,83.5,87,90.5,94,97.5}{\draw[thick,gray!40!black] (\s,5.25) -- (\s,5.65);}
\draw[gray!70!black] (0,5.25) -- (105,5.25);
\node[left,font=\small] at (-1,5.45) {المدخل};
\node[above,font=\scriptsize,gray!40!black] at (13.5,5.65) {نادرًا};
\node[above,font=\scriptsize,gray!40!black] at (21.5,6.0) {زوج};
\node[above,font=\scriptsize,gray!40!black] at (75,5.65) {كثيرًا};
% axes and thresholds
\foreach \y/\th in {2.7/2.5,0/1.5}{
  \draw[gray!70!black] (0,\y) -- (106,\y) node[right,black] {\small $t$};
  \draw[densely dashed,gray!60!black] (0,\y+0.6*\th) -- (104,\y+0.6*\th) node[right,black,font=\small] {$\theta$};
}
\node[anchor=south west,font=\small] at (0,4.3) {$\tau=10\ms$: مكامل};
\node[anchor=south east,font=\small] at (104,1.0) {$\tau=2\ms$: كاشف تزامن};
% integrator: tau=10 ms, theta=2.5w
\begin{scope}[yshift=2.7cm]
\foreach \a/\b/\v in {0/6/0.000,6/21/1.000,21/22/1.223,22/38/2.107,38/52/1.425,52/55.5/1.351,55.5/59/1.952,59/62.5/2.376,62.5/66/0.000,66/69.5/1.000,69.5/73/1.705,73/76.5/2.201,76.5/80/0.000,80/83.5/1.000,83.5/87/1.705,87/90.5/2.201,90.5/94/0.000,94/97.5/1.000,97.5/104/1.705}{\draw[very thick,blue!60!black] plot[domain=\a:\b,samples=14] (\x,{0.6*\v*exp(-(\x-\a)/10)});}
\foreach \s/\p/\q in {6/0.000/1.000,21/0.223/1.223,22/1.107/2.107,38/0.425/1.425,52/0.351/1.351,55.5/0.952/1.952,59/1.376/2.376,62.5/1.674/2.500,62.5/2.500/0.000,66/0.000/1.000,69.5/0.705/1.705,73/1.201/2.201,76.5/1.551/2.500,76.5/2.500/0.000,80/0.000/1.000,83.5/0.705/1.705,87/1.201/2.201,90.5/1.551/2.500,90.5/2.500/0.000,94/0.000/1.000,97.5/0.705/1.705}{\draw[very thick,blue!60!black] (\s,0.6*\p) -- (\s,0.6*\q);}
\foreach \s in {62.5,76.5,90.5}{\draw[->,thick,red!70!black] (\s,1.58) -- (\s,1.95);}
\end{scope}
% coincidence: tau=2 ms, theta=1.5w
\begin{scope}[yshift=0cm]
\foreach \a/\b/\v in {0/6/0.000,6/21/1.000,21/22/1.001,22/38/0.000,38/52/1.000,52/55.5/1.001,55.5/59/1.174,59/62.5/1.204,62.5/66/1.209,66/69.5/1.210,69.5/73/1.210,73/76.5/1.210,76.5/80/1.210,80/83.5/1.210,83.5/87/1.210,87/90.5/1.210,90.5/94/1.210,94/97.5/1.210,97.5/104/1.210}{\draw[very thick,blue!60!black] plot[domain=\a:\b,samples=14] (\x,{0.6*\v*exp(-(\x-\a)/2)});}
\foreach \s/\p/\q in {6/0.000/1.000,21/0.001/1.001,22/0.607/1.500,22/1.500/0.000,38/0.000/1.000,52/0.001/1.001,55.5/0.174/1.174,59/0.204/1.204,62.5/0.209/1.209,66/0.210/1.210,69.5/0.210/1.210,73/0.210/1.210,76.5/0.210/1.210,80/0.210/1.210,83.5/0.210/1.210,87/0.210/1.210,90.5/0.210/1.210,94/0.210/1.210,97.5/0.210/1.210}{\draw[very thick,blue!60!black] (\s,0.6*\p) -- (\s,0.6*\q);}
\foreach \s in {22}{\draw[->,thick,red!70!black] (\s,0.98) -- (\s,1.35);}
\end{scope}
\foreach \x in {0,50,100}{\draw[gray!60!black] (\x,-0.06) -- (\x,0.06) node[below,font=\scriptsize] {$\x\,\text{ms}$};}
\end{tikzpicture}
\caption{المدخل نفسه، ومستقبِلان مختلفان. كل نبضة ترفع الجهد بمقدار $w$، ثم ينساب نازلًا كما في نموذج العصبون~\eqref{eq:glutneuron}؛ والأسهم الحمراء نبضات المستقبِل. المستقبِل البطيء ($\tau=10\ms$، $\theta=2.5w$) ينطلق حيث النبضات \emph{كثيرة}، ولا يلاحظ الزوج. والسريع ($\tau=2\ms$، $\theta=1.5w$) لا ينطلق إلا للزوج الواقع ضمن النافذة~\eqref{eq:window}، ويُفلت التدفق المتواتر غير المتزامن.}
\label{fig:reader}
\end{figure}

تُظهر القياسات أن موصلية الغشاء في القشرة النشطة تزداد عدة مرات مقارنة بالراحة~\cite{Destexhe2003}. إذن $\tau$ هناك أقصر، وعصبونات القشرة في نمط العمل أقرب إلى كواشف التزامن منها إلى المكاملات.

\begin{keyblock}
\begin{proposition}[المستقبِل يحدد الشيفرة]
\label{prop:reader}
متتالية النبضات نفسها تحمل للمستقبِل البطيء ترددًا، وللسريع أزمنة، وللحجم الذي قُذف فيه الناقل العصبي مستوى تركيز منعّمًا على مدى ثوانٍ (الفصل~\ref{sec:serocomputer}). أي شيفرة تُستعمل لا يحدده المرسِل، بل الثابت الزمني للمستقبِل، والتثبيط هو الذي يضبطه.
\end{proposition}
\end{keyblock}
\section{النقاط والشُّرَط: المشبك مرشِّحًا}

حتى الآن كان المشبك عندنا سلكًا ذا وزن. والحقيقة أن استجابته تتوقف على النبضات السابقة، وهذا يعطي لغة العصبونات رمزًا ثانيًا.

\emph{النضوب: المشبك ينقل التغيرات.} كل قذفة تستهلك حصة $U$ من مخزون الناقل العصبي الجاهز للقذف~\cite{Tsodyks1997}، ويُستعاد المخزون بثابت زمني $\tau_{\text{rec}}$ من رتبة مئات من أجزاء الألف من الثانية. عند التردد $r$ تستقر سعة الاستجابة للنبضة تقريبًا عند المستوى $A(r)=A_0/(1+U r\tau_{\text{rec}})$، حيث $A_0$ استجابة المشبك المرتاح. والتيار الذي يتلقاه المستقبِل في وحدة الزمن يساوي
\begin{empheq}[box=\keybox]{equation}
r\,A(r) \;=\; \frac{A_0\,r}{1+U r\,\tau_{\text{rec}}}
\;\xrightarrow[r\to\infty]{}\; \frac{A_0}{U\tau_{\text{rec}}} .
\label{eq:depression}
\end{empheq}
عند $U=0.5$ و$\tau_{\text{rec}}=0.8\,\text{s}$ يبدأ الإشباع منذ $1/(U\tau_{\text{rec}})=2.5$ نبضة في الثانية. فإذا ارتفع التردد من خمس إلى عشرين، لم يزد التيار المستقر إلا بالثلث. لكن النبضات الأولى على التردد الجديد تصل والمخزون ما زال على حاله، فيقفز التيار لحظةً أربعة أضعاف، ثم يهبط خلال $\tau_{\text{rec}}/(1+Ur\tau_{\text{rec}})\approx90\ms$ (الشكل~\ref{fig:depression}).

هذا المشبك لا ينقل التردد بل \emph{تغيّره}، وفي جوهر الأمر التغير النسبي $\Delta r/r$~\cite{Abbott1997depression}. وبالنسبة لمهندس الإلكترونيات هذا مرشِّح تمرير عالٍ موضوع في السلك مباشرة.

\begin{figure}[t]
\centering
\begin{tikzpicture}[>=Stealth,x=12.5cm,y=1cm]
% input rate
\draw[->,gray!70!black] (0,2.6) -- (0.9,2.6) node[right,black] {\small $t$};
\draw[->,gray!70!black] (0,2.6) -- (0,4.3) node[above,black] {\small $r$};
\draw[very thick,blue!60!black] (0,2.9) -- (0.2,2.9) -- (0.2,3.8) -- (0.8,3.8);
\node[left,font=\scriptsize] at (0,2.9) {$5$};
\node[left,font=\scriptsize] at (0,3.8) {$20$};
\node[right,font=\small,blue!60!black] at (0.4,4.1) {التردد عند مدخل المشبك};
% output current
\draw[->,gray!70!black] (0,0) -- (0.9,0) node[right,black] {\small $t$};
\draw[->,gray!70!black] (0,0) -- (0,2.45) node[left,black] {\small $rA$};
\draw[densely dashed,gray!60!black] (0,0.667) -- (0.8,0.667);
\draw[very thick,red!70!black] (0,0.5) -- (0.2,0.5) -- (0.2,2.0)
   plot[domain=0.2:0.8,samples=60] (\x,{0.667+(2.0-0.667)*exp(-(\x-0.2)/0.089)});
\node[left,font=\scriptsize] at (0,0.5) {$1.7$};
\node[left,font=\scriptsize] at (0,2.0) {$6.7$};
\node[right,font=\scriptsize] at (0.8,0.667) {$2.2$};
\node[right,font=\small,red!70!black] at (0.28,1.6) {التيار إلى المستقبِل: قفزة ثم هبوط خلال $\approx90\ms$};
\draw[gray!60!black] (0.2,-0.08) -- (0.2,0.08); \node[below,font=\scriptsize] at (0.2,0) {$0.2\,\text{s}$};
\draw[gray!60!black] (0.8,-0.08) -- (0.8,0.08); \node[below,font=\scriptsize] at (0.8,0) {$0.8\,\text{s}$};
\end{tikzpicture}
\caption{مشبك ناضب بحسب الصيغة~\eqref{eq:depression} عند $U=0.5$ و$\tau_{\text{rec}}=0.8\,\text{s}$؛ التيار بوحدات $A_0$ في الثانية، والخط المتقطع هو التيار المستقر: لم يرتفع إلا من $1.7$ إلى $2.2$، أما لحظة القفزة فإلى $6.7$. يبلغ المشبك المستقبِلَ أن التردد تغيّر، لا ما هو.}
\label{fig:depression}
\end{figure}

\emph{التيسير: الرشقة شَرطة.} في مشابك أخرى احتمال القذف صغير، لكنه يزداد بعد كل نبضة لعشرات من أجزاء الألف من الثانية. النبضة المفردة عبر مشبك كهذا تضيع في الغالب. أما \emph{الرشقة}، أي بضع نبضات متتالية تفصل بينها بضعة أجزاء من الألف من الثانية، فتمر شبه مؤكدة. حتى من دون تيسير، احتمال أن تضيع نبضات الرشقة الـ$k$ كلها يساوي
\begin{empheq}[box=\keybox]{equation}
P_{\text{loss}} \;=\; (1-p)^k ,
\label{eq:burst}
\end{empheq}
وعند $p=0.2$ هذا $0.8$ للنبضة المفردة و$0.41$ لرشقة من أربع؛ والتيسير يخفضه أكثر. هكذا يظهر لدى العصبون رمز ثانٍ: النبضة المفردة نقطة، والرشقة شَرطة. المشبك الناضب أفضل ما ينقل النبضة الأولى بعد توقف؛ والمشبك التيسيري يمرر الرشقات ويخمد النقاط المفردة. أما أن الرشقات تنتقل بموثوقية أكبر فمُقاس؛ وأما أن الدماغ يستعملها فعلًا رمزًا مستقلًا ففرضية معقولة~\cite{Lisman1997}.

\section{كيف تتكلم عصبونات \nt{GABA}}

أكثرها عددًا في القشرة هي السريعة~\cite{Tremblay2016}. نبضاتها أقصر من نبضات عصبونات \nt{GLUT}، وتستطيع أن تطلقها بانتظام ولمدة طويلة بتردد مئات في الثانية، بلا تعب يُذكر. ومشابكها أوثق: فالوصلة مع كل مستقبِل مؤلفة من مواضع قذف كثيرة، ولا تكاد النبضة تضيع أبدًا. وتستقر مخارجها قرب الموضع الذي تولد فيه نبضة المستقبِل، فهي لا تتحكم في كيفية جمع المستقبِل لمداخله، بل في هل ينطلق ومتى.

وهي نفسها تتلقى الإثارة من عصبونات \nt{GLUT} متجاورة كثيرة وتُبلغ عن الفعالية الكلية من حولها، وهذا بالضبط ما تحتاجه القسمة في الفصل~\ref{sec:gaba}. وأخيرًا، عصبونات \nt{GABA} السريعة متصلة بعضها ببعض وتنطلق معًا بسهولة؛ وهي بالذات التي تحدد الإيقاع المحلي الذي تحدثنا عنه أعلاه~\cite{Cardin2009}.

بلغة شفرة مورس، عصبونات \nt{GABA} ليست مسؤولة عن الحروف، بل عن \emph{الفواصل}. تقطّع التدفق المتصل لنبضات \nt{GLUT} إلى نوافذ يمكن للمستقبِل أن ينطلق داخلها، وتضيّق نوافذ التزامن، وتخفت التدفق كله حين يعلو أكثر مما ينبغي.

\begin{keyblock}
\begin{proposition}[تقسيم الأدوار]
\label{prop:punctuation}
عصبونات \nt{GLUT} تحمل مضمون الرسالة: \emph{ماذا} و\emph{كم}. وعصبونات \nt{GABA} السريعة تحمل علاماتها: \emph{متى} يجوز الكلام و\emph{بأي علوّ}؛ وصنف آخر، بتثبيطه للمثبِّطات، يقرر \emph{لمن} تُفتح البوابة. بعض أجزاء هذا التقسيم مُقاس؛ أما بوصفه قاعدة عامة ففرضية عمل.
\end{proposition}
\end{keyblock}

وثمة عصبونات \nt{GABA} أخرى أبطأ، تستقر مخارجها على مداخل منفردة للمستقبِل. تعمل كالحظر~\eqref{eq:veto} في الفصل~\ref{sec:gaba}: تغلق تدفقًا واحدًا من نبضات \nt{GLUT} دون المساس بغيره.

تقسيم عصبونات \nt{GABA} إلى سريعة وبطيئة صورة قديمة، وهي ناقصة. فاليوم يُميَّز في القشرة ثلاثة أصناف كبرى على الأقل~\cite{Tremblay2016}، والثالث مصمم على نحو غير متوقع: لا يثبط عصبونات \nt{GLUT}، بل عصبونات \nt{GABA} أخرى، ولا سيما تلك التي تغلق المداخل~\cite{Pfeffer2013}. تثبيط التثبيط نفي مزدوج. عصبون كهذا \emph{يفتح} البوابة دون أن يرسل إلى المستقبِل نبضة مثيرة واحدة. تشغّله إشارات من مناطق أخرى في القشرة ومن قنوات البثّ، فيعمل مدخلَ تمكين لجزء كامل من الشبكة: ما دام نشطًا فالحظور مرفوعة، والتدفق يمر~\cite{Pi2013}. وفي الملحق~\ref{sec:othermediators} سُميت هذه الحيلة رفع التثبيط.
\section{لغة مقتصدة}

آخر ما يُعرف بثبات عن لغة العصبونات أنها مكلفة. النبضة، مع عمل المشابك الذي تستدعيه، تكلف نحو $10^9$ جزيء من ATP (تقدير لقشرة القوارض~\cite{Attwell2001})، والجزيء الواحد يعطي نحو $10^{-19}$ جول. إذن تكلف النبضة من رتبة $E_{\text{spk}}\sim10^{-10}$ جول. يستهلك الدماغ كله نحو $20$ واط، وإذا ذهب نصفها إلى نقل الإشارات، $P\sim10$ واط، فالتردد المتوسط للعصبون الواحد
\begin{empheq}[box=\keybox]{equation}
\bar r \;\approx\; \frac{P}{N\,E_{\text{spk}}}
\;\sim\; \frac{10\ \text{W}}{8.6\cdot10^{10}\cdot10^{-10}\ \text{J}}
\;\approx\; 1\ \text{spike/s}.
\label{eq:energy}
\end{empheq}
والتقديرات الأكثر تفصيلًا للقشرة تعطي أقل من ذلك~\cite{Lennie2003}. شيفرة الدماغ مضطرة إلى أن تكون \emph{متناثرة}: لا يستطيع أن يعمل بسرعة إلا جزء صغير من العصبونات.

ومن~\eqref{eq:spikeentropy} يتبين أن هذا ليس سيئًا إلى هذا الحد. بدقة $1\ms$ لا تحمل نبضة عصبون يعمل بتردد $100$ في الثانية أكثر من خمسة بتات، أما نبضة عصبون يعمل مرة في الثانية فتحمل حتى أحد عشر. وإذا كان الدفع على النبضات، فالأربح أن تتكلم نادرًا. والقشرة فعلًا تقايض الدقة بالطاقة: عند شح الغذاء تحافظ العصبونات على تردد نبضاتها، لكنها تنفق أقل على التيارات المشبكية، فتصبح استجاباتها أقل دقة~\cite{Padamsey2022}.

في شفرة مورس الحروف الشائعة قصيرة: أشيع حرف في النص الإنجليزي، \foreignlanguage{english}{E}، نقطة واحدة. وشيفرة العصبونات، بحسب ما هو معروف، تتبع القاعدة نفسها في صورة أخرى: المتوقع يكلف نبضات قليلة~\cite{Barlow1961}. العصبون يخفض تردده تدريجيًا مع المنبه الثابت، ومستشعرات الفصل~\ref{sec:glutcomputer} تستجيب للتغيرات لا للمستويات، وعصبونات \nt{DOPA} في الفصل~\ref{sec:neurontypes} لا تبلغ إلا عن خطأ التنبؤ بالمكافأة.

\begin{keyblock}
\begin{proposition}[الاقتصاد]
\label{prop:economy}
تحدّ الطاقة التردد المتوسط للعصبون بقيمة من رتبة نبضة واحدة في الثانية. لذلك فلغة عصبونات \nt{GLUT} متناثرة، وتنفق نبضاتها على الأخبار: على ما تغيّر أو لم يُتنبأ به. الحد الطاقي مُقاس؛ أما أن شيفرة الدماغ مضبوطة على أكبر عدد من البتات لكل جول ففرضية حسنة الأسس.
\end{proposition}
\end{keyblock}

\section{الخلاصة}

\begin{table}[t]
\centering
\footnotesize
\setlength{\tabcolsep}{4pt}
\renewcommand{\arraystretch}{1.15}
\begin{tabular}{>{\raggedright}p{2.6cm}>{\raggedright}p{2.3cm}>{\raggedright}p{3.0cm}>{\raggedright}p{2.0cm}>{\raggedright\arraybackslash}p{3.6cm}}
\toprule
طريقة الكتابة & ماذا تحمل & من يقرؤها & زمن الرسالة & ما هو معروف \\
\midrule
رقم العصبون المنطلق (الترميز بالموضع) & أي قيمة & أي مستقبِل؛ كواشف التزامن & تأخير واحد & ثابت في العصبونات الحسية وفي تحديد اتجاه الصوت \\
تردد عصبون واحد & المقدار & مكامل بـ$\tau$ كبير؛ الحجم (الفصل~\ref{sec:serocomputer}) & من مئات الأجزاء من الألف من الثانية إلى ثوانٍ & ثابت؛ وبطيء جدًا لمرحلة معالجة من $10\text{--}20\ms$ \\
تردد مجموعة & المقدار & عصبون بآلاف المداخل & $10\text{--}50\ms$ & ثابت؛ والربح تحده الارتباطات~\eqref{eq:correlated} \\
زمن النبضة الأولى، الترتيب & المقدار، الترتيب & كواشف التزامن ذات التأخيرات & تأخير واحد & ثابت عند مدخل الجهاز العصبي؛ وفي القشرة موضع نقاش \\
الطور في الإيقاع المحلي & المقدار، الموضع & عصبونات ذات إيقاع مشترك & دورة $12\text{--}50\ms$ وأبطأ & مرصود في مناطق بعينها؛ ودوره العام فرضية \\
الرشقة (”الشَّرطة“) & ”مهم“: إيصال موثوق & المشبك التيسيري & $\sim10\ms$ & الموثوقية مقيسة؛ والرمز المستقل فرضية \\
تغيّر التردد & خبر & المشبك الناضب & $\sim0.1\,\text{s}$ & ثابت \\
نبضات عصبونات \nt{GABA} السريعة & متى وبأي علوّ & كل عصبونات \nt{GLUT} المجاورة & $1\text{--}30\ms$ & الإيقاع والقسمة مقيسان؛ و”علامات الترقيم“ قاعدةً فرضية \\
\bottomrule
\end{tabular}
\caption{الطرق التي تكتب بها عصبونات \nt{GLUT} و\nt{GABA} الرسائل بالنبضات. العمود الأيسر يفصل المقيس عن المفترض.}
\label{tab:code}
\end{table}

يجمع الجدول~\ref{tab:code} ما نعرفه؛ ويظهر منه أيضًا ما لا نعرفه. لا نعرف إلى أي حد تستعمل القشرة التوقيت الدقيق للنبضات، لا ترددات المجموعات فقط. ولا نعرف هل للعصبونات ”كلمات“، أي توليفات مستقرة من نبضات عصبونات كثيرة تُقرأ كلًّا واحدًا. ولا نعرف هل اللغة واحدة في أجزاء الدماغ المختلفة. يمكن تعلم شفرة مورس في أمسية، لأن جدولها معروف. أما جدول لغة العصبونات فما زال ينتظر من يضعه، والأرجح أن من سيضعه مهندسو الاتصالات.

\section{تمارين}

\begin{exercise}[\lvlA]
\label{ex:04:1}
\emph{السعة بالأرقام.} $\Delta t=1\ms$. (أ)~كم بتًا لكل جول يعطي عصبون تردده $1$ نبضة في الثانية، بكلفة النبضة من~\eqref{eq:energy}؟ (ب)~عند $r=10$، كم مرة أقل من الحد~\eqref{eq:spikeentropy} يستخرج مستقبِل لا يفعل سوى عدّ النبضات في الثانية؟
\end{exercise}

\begin{exercise}[\lvlB]
\label{ex:04:2}
\emph{التردد الأمثل.} ينفق العصبون قدرة $P_0+rE_{\text{spk}}$، حيث $P_0$ هو النصف الآخر من $20\,\text{W}$ مقسومًا على كل العصبونات، و$E_{\text{spk}}=10^{-10}\,\text{J}$. عند أي تردد $r$ يبلغ عدد البتات لكل جول $R_{\max}(r)/(P_0+rE_{\text{spk}})$ بحسب~\eqref{eq:spikeentropy} مع $\Delta t=1\ms$ أقصاه؟
\end{exercise}

\begin{exercise}[\lvlB]
\label{ex:04:3}
\emph{أي ارتباط ضار.} $N=1000$ عصبون، التباين $\sigma^2=1$، والارتباط الثنائي $c=0.1$. احسب معلومات فيشر $\mathcal I=f'^{\mathsf T}\Sigma^{-1}f'$ ($\Sigma$ مصفوفة التغاير) لميول متساوية $f'=\mathbf 1$ ولميول $\pm1$ مع $\mathbf 1^{\mathsf T}f'=0$. كم مرة تختلفان؟
\end{exercise}

\begin{exercise}[\lvlB]
\label{ex:04:4}
\emph{نمذجة: كاشف التزامن.} العصبون~\eqref{eq:glutneuron} يتلقى $1000$ مدخل بواسوني بـ$5$ نبضات في الثانية لكل منها، $w=1$؛ والعتبة $\theta$ بحيث يعبرها الجهد الحر مرة في الثانية. أوجد بالنمذجة كم مدخلًا متزامنًا $m$ يُطلق العصبون باحتمال $1/2$ عند $\tau=10$ و$2\ms$.
\end{exercise}

\begin{exercise}[\lvlB]
\label{ex:04:5}
\emph{تصميم: كاشف التغيرات النسبية.} اختر مشبكًا~\eqref{eq:depression}: التيار المستقر عند $40$ نبضة في الثانية لا يزيد بأكثر من $10\,\%$ عليه عند $10$، وبعد القفزة إلى $40$ يبلغ التيار المستوى الجديد بثابت زمني لا يتجاوز $50\ms$. ما أصغر $\tau_{\text{rec}}$ وعند أي $U$؟
\end{exercise}

\begin{exercise}[\lvlA]
\label{ex:04:6}
\emph{زمن النبضة الأولى والرشقات.} (أ)~بحسب~\eqref{eq:latency} أوجد المدخل الذي تأتي عنده النبضة الأولى بعد ثابت زمني واحد بالضبط. علامَ يتوقف؟ (ب)~عند احتمال قذف $p=0.2$، كم نبضة تلزم في الرشقة ليكون احتمال ضياعها كلها أقل من $5\,\%$؟
\end{exercise}

\begin{exercise}[\lvlC]
\label{ex:04:7}
\emph{بحث: كم بتًا في الترتيب المكاني.} العصبون خانات مستقلة~\eqref{eq:spikeentropy}، $r=10$، $\Delta t=1\ms$. (أ)~فكّك إنتروبيا ”كلمة“ من $20$ خانة إلى إنتروبيا عدد النبضات والباقي. (ب)~نمذج تقدير الإنتروبيا من ترددات الكلمات بـ$N=30$ و$10^4$ تسجيل. بكم يكون التقدير منخفضًا؟
\end{exercise}
